% Trigonometrie · Sinussatz · Prüfungs-Fokus MSA – bank/trigonometrie/e4.jsonl (zusammenbau v0.6, Prüfungs-Fokus)
\einheitenkopf[e4]{Einheit 4 · Sinussatz}
\zweigzeile{ab Kl. 10 (am Gymnasium ab Kl. 9) · P10 ×9}
\setcounter{aufgabe}{0}

% Anlauf (ohne Prüfkennung)
\begin{aufgabe}{Sinussatz – Anlauf}
\begin{teile}
\teil Im Dreieck $ABC$ ist bei $B$ ein rechter Winkel markiert, $a$ und $\alpha$ sind gegeben. Womit rechnest du $b$? \\ \kreuz{sin, cos, tan} \kreuz{Sinussatz}
\teil Dreieck $ABC$: $a = 7$ cm, $\alpha = 68^\circ$, $\beta = 51^\circ$. Wie lang ist $b$? b ≈ \leerfeld[cm]
\teil Dreieck $PQR$: $PQ = 8$ cm, Winkel bei $R$: $69^\circ$, Winkel bei $Q$: $53^\circ$. Wie lang ist $PR$? PR ≈ \leerfeld[cm]
\end{teile}
\end{aufgabe}

\begin{aufgabe}{Sinussatz – Prüfungsaufgaben}
\begin{teile}
\steil Im Drachen $ABCD$ gilt im Dreieck $ABD$: $AD = 4{,}6$ cm, der Winkel bei $A$ ist $118^\circ$, der Winkel bei $B$ ist $33^\circ$. Berechne die Diagonale $BD$. \hfill (P15*) \\ BD ≈ \leerfeld[cm]
\steil Wanderung: Hütte $H$, Gipfel $G$, See $S$. $HG = 3\,200$ m, der Winkel bei $H$ ist $57^\circ$, der Winkel bei $S$ ist $72^\circ$. Wie lang ist der Weg vom Gipfel zum See und weiter zum Parkplatz, der $1\,800$ m hinter dem See liegt? \hfill (P14*) \\ \leerfeld[km]
\end{teile}
\end{aufgabe}

\begin{aufgabe}{Sinussatz – Prüfungsaufgaben – weiter}
\begin{teile}
\steil Eine Seilbahn führt von $A$ über die Stütze $B$ nach $C$. $AB = 520$ m, der Winkel bei $A$ ist $31^\circ$, der Winkel $ABC$ ist $113^\circ$. Berechne die Länge der Strecke $BC$. \hfill (P24*) \\ BC ≈ \leerfeld[m]
\steil Eine Seilbahn führt von $A$ über die Stütze $B$ nach $C$. $AB = 290$ m, der Winkel bei $A$ ist $27^\circ$, der Winkel $ABC$ ist $121^\circ$. Berechne die Länge der Strecke $BC$. \hfill (P24*) \\ BC ≈ \leerfeld[m]
\steil Eine Rampe $y$ führt auf eine Treppe. Im Dreieck aus Rampe, waagerechter Strecke ($220$ cm) und der Verbindung Treppenfuß–obere Stufenkante liegt am Rampenfuß der Winkel $6^\circ$ und am Treppenfuß der Winkel $137^\circ$ (gegenüber $y$). Berechne $y$. \hfill (P25*) \\ y ≈ \leerfeld[cm]
\steil Eine Rampe $y$ führt auf eine Treppe. Im Dreieck aus Rampe, waagerechter Strecke ($190$ cm) und der Verbindung Treppenfuß–obere Stufenkante liegt am Rampenfuß der Winkel $4^\circ$ und am Treppenfuß der Winkel $148^\circ$ (gegenüber $y$). Berechne $y$. \hfill (P25*) \\ y ≈ \leerfeld[cm]
\steil Von einer Plattform aus werden zwei Sichtlinien gepeilt. Im Dreieck $ACD$ ist $CD = 9$ m, der Winkel bei $A$ zwischen den Sichtlinien beträgt $4^\circ$, der Winkel bei $D$ $62^\circ$. Zeige, dass $AC \approx 114$ m ist. \hfill (P18*)
\steil Von einer Plattform aus werden zwei Sichtlinien gepeilt. Im Dreieck $ACD$ ist $CD = 8$ m, der Winkel bei $A$ zwischen den Sichtlinien beträgt $6^\circ$, der Winkel bei $D$ $51^\circ$. Zeige, dass $AC \approx 59{,}5$ m ist. \hfill (P18*)
\steil Ein Grundstück $ABCD$ hat bei $B$ einen rechten Winkel. $AC = 7$ m, $\angle BAC = 59^\circ$, $\angle BCD = 88^\circ$, $\angle ADC = 78^\circ$. Berechne den Abstand $AD$. \hfill (P19*) \\ AD ≈ \leerfeld[m]
\steil Ein Grundstück $ABCD$ hat bei $B$ einen rechten Winkel. $AC = 10$ m, $\angle BAC = 53^\circ$, $\angle BCD = 96^\circ$, $\angle ADC = 69^\circ$. Berechne den Abstand $AD$. \hfill (P19*) \\ AD ≈ \leerfeld[m]
\steil Dreieck $DEF$: $FE = 9$ m, Winkel bei $F$ $118^\circ$, Winkel bei $E$ $23^\circ$. $P$ liegt auf $DE$ mit $PE = 8$ m. Wie lang ist $DP$? \hfill (P20*) \\ DP ≈ \leerfeld[m]
\steil Dreieck $DEF$: $FE = 16$ m, Winkel bei $F$ $107^\circ$, Winkel bei $E$ $27^\circ$. $P$ liegt auf $DE$ mit $PE = 17$ m. Wie lang ist $DP$? \hfill (P20*) \\ DP ≈ \leerfeld[m]
\steil Viereck $ABCD$ mit rechtem Winkel bei $A$ und der Diagonale $BD$. $\angle ADB = 48^\circ$, $\angle BDC = 57^\circ$, $\angle ABC = 104^\circ$, $DC = 1\,300$ m. Stimmt die Aussage „$BC$ ist genauso lang wie $DC$“? \hfill (P21*)
\steil Viereck $ABCD$ mit rechtem Winkel bei $A$ und der Diagonale $BD$. $\angle ADB = 51^\circ$, $\angle BDC = 63^\circ$, $\angle ABC = 112^\circ$, $DC = 1\,760$ m. Stimmt die Aussage „$BC$ ist genauso lang wie $DC$“? \hfill (P21*)
\end{teile}
\end{aufgabe}

\medskip
\uebersichtskasten{Sinussatz: In jedem Dreieck ist das Verhältnis aus einer Seite und dem Sinus ihres Gegenwinkels für alle drei Seiten gleich – a/sin α = b/sin β = c/sin γ. Man braucht ein vollständiges Paar (Seite und Gegenwinkel) und von einem zweiten Paar eine Größe. \\ b = 8 cm, β = 40°, α = 75°: a = 8 · sin 75° : sin 40° ≈ 12,0 cm \\ Fehlt der Winkel gegenüber der gesuchten Seite, liefert ihn die Winkelsumme: γ = 180° − 75° − 40° = 65°. \\ Ein stumpfer Winkel bleibt, wie er ist – der Taschenrechner rechnet sin 120° direkt; in die Winkelsumme gehört der stumpfe Winkel selbst, nicht sein Nebenwinkel.}
